\chapter{Ownership, capacity and atomic accounting}\label{ch:accounting}
The most important unimplemented boundary in the new programme is not a
missing sum. It is the authority to interpret a particular sum as an owner's
capacity change. Software can enforce an adopted rule, but it cannot obtain
that rule from the fact that several numbers add correctly.

\section{A child action is not its owner}
A request can be made by a clinic, performed by a carrier and supplied by
a depot. The action record may identify all three. None of those labels
automatically defines who owns its path attribution. The prospective
programme needs an explicit, versioned $\operatorname{owner}(a)$ map.

Once each child has one declared owner, aggregation is straightforward:
$\Delta B_i=\sum_{a:\operatorname{owner}(a)=i}R_a$. Closure follows
by regrouping the finite sum. The difficult question is not whether the
sum can be implemented, but which policy it represents and what claims
may be made about it.

\section{Same-event netting is a policy}
Suppose one owner has child lines $+5$ and $-3$ and starts at zero.
Owner-level same-event netting produces $+2$ and passes the proposed
affordability rule. A policy requiring each negative child to be prefunded
would refuse the negative child. A global-pooling policy could do something
else again. These are distinct stochastic processes.

Likewise, a group with owner changes $(+5,-3)$ is not affordable from
two zero balances merely because its total is positive. One owner's surplus
does not become another's funding under the no-transfer proposal. The
implementation must not substitute a single aggregate-balance check for
owner-wise checks.

\section{All-or-nothing effects}
If a selected group is accepted under a joint contract, physical and account
updates must refer to that same group. A partial credit followed by a failed
debit can create spendable capacity without the corresponding authorized
event. A physical update followed by a lost receipt can make replay execute
the same action twice.

A later execution design must therefore define commit boundaries, durable
identities and recovery behavior across the linked records. ``Atomic'' here
means an all-or-nothing committed event under that contract; it does not
mean hiding a multi-carrier physical chain inside one mega-action.

\section{Audit is not a spendable pool}
The external audit $J$ records signed forcing effects. It is not a wallet,
loan facility or reward for the next actor. In the cost-free proposal,
$V+\sum_iB_i-J=K$. A software routine must be able to check that
identity without making $J$ available to the selector as spending capacity.

The invariant alone cannot identify which record is wrong. An erroneous
physical update and matching erroneous credit can cancel. Independent
conservation, action-identity and owner-level checks are necessary in
addition to the aggregate account residual.

\section{Existing settlement machinery has a narrower meaning}
The source-locked framework includes a repaired share-plus-residual closure
validator on another lineage \cite{framework}. That is an implemented
validation mechanism, not an adopted Gaussian owner rule. Its historical
amendment's authoring-stage label is not the whole later implementation
history; neither should later implementation be mistaken for scientific
approval of the new capacity mechanism.

\section*{Chapter recap}
Capacity requires a declared ownership and update policy, exact owner-wise
affordability, linked event identity and durable all-or-nothing effects.
The book specifies the boundary; it does not claim that runtime exists.
